\begin{abstract}
Model merging has emerged as a powerful paradigm for combining multiple task-specific expert neural networks into a single multi-task model without additional training. However, the practical deployment of merged models requires compression, such as post-training quantization (PTQ), which introduces significant noise and often severely degrades performance. In this paper, we present \textbf{FlatQ-Merge} (Flatness-Aware Quantization-Aware Model Merging), an exhaustive empirical study of the relationship between the loss landscape flatness of expert models and their subsequent resilience to weight quantization and test-time coefficient optimization.

We pre-train task-specific experts across 5 different Sharpness-Aware Minimization (SAM) perturbation radii ($\rho \in \{0.0, 0.01, 0.05, 0.1, 0.2\}$) and evaluate their merged, quantized forms under 8-bit and extreme 4-bit precision across 3 independent random seeds. Our extensive multi-axial grid sweeps on a Vision Transformer (\texttt{vit\_tiny\_patch16\_224}) backbone across MNIST, FashionMNIST, CIFAR-10, and SVHN uncover a powerful, precision-dependent \emph{Flatness-Robustness Synergy}. Specifically, we find that the flatness-robustness synergy is highly precision-dependent: while standard 8-bit formats are highly robust even with SGD-trained experts, pre-training experts with an optimal SAM radius ($\rho=0.05$) yields a massive \textbf{+7.44\%} absolute multi-task accuracy improvement under extreme 4-bit quantization. We also identify a critical non-linear \emph{Over-Perturbation Threshold} at $\rho \ge 0.1$ where performance collapses.

Remarkably, we discover that merging flat experts naively with static uniform weights ($\rho=0.05$, NaiveUniform) outperforms performing sophisticated test-time adaptation on sharp experts ($\rho=0.0$, FlatQ-Merge) by a massive \textbf{+6.03\%} absolute accuracy. This demonstrates that pre-merging expert geometry is significantly more critical to low-precision merging success than the complexity of the downstream test-time adaptation algorithm itself.
\end{abstract}
